\documentclass[11pt, a4paper, oneside]{article}
\usepackage[ngerman]{babel}
\usepackage{worksheet}

\begin{document}
	\author{L. Bung}
	\title{Logarithmen}
	\subject{Mathematik}
	\class{FHR1}
	\maketitle
	
	\singletask{Einfache Exponentialgleichungen lösen}
	
	Bestimmen Sie die Lösung der Exponentialgleichungen.
	
	\begin{multicols}{3}
		\begin{enumerate}[label=\alph*)]
			\item $2^x=32$
			\item $2\cdot3^x=54$
			\item $3\cdot2^{2x}=48$
		\end{enumerate}
	\end{multicols}
	
	\checkered[6cm]
	
	\begin{hint}{Der Logarithmus}
		Der \textbf{Logarithmus} zur Basis $q$ -- geschrieben $\log_q(x)$ -- ist die Umkehrfunktion zur Exponentialfunktion $q^x$ mit der Basis $q$.
		Das bedeutet: $\log_q(q^x) = x$ und $q^{\log_q(x)}=x$.
		
		\vspace{.5cm}
		Die Logarithmen mit häufig verwendeten Basen haben besondere Schreibweisen:
		\begin{multicols}{2}
			Zehnerlogarithmus (Basis 10): $\log_{10} \widehat{=} \lg$
			
			Natürlicher Logarithmus (Basis $e$): $\log_e \widehat{=} \ln$
		\end{multicols}
		
		Exponentialgleichungen, die nicht so einfach im Kopf zu rechnen sind, lassen sich mit dem Logarithmus lösen:
		
		$$\begin{aligned}
		2^x &= 12 &&|\log_2()\\
		x &= \log_2(12)\\
		x &\approx 3,58
		\end{aligned}$$
	\end{hint}
	
	\singletask{Exponentialgleichungen finden}
	
	Finden Sie Gleichungen, die die angegebene Lösung haben.
	\emph{Beispiel}: Zur Lösung $\ln(2)$ wäre eine mögliche Gleichung $e^x=2$.
	
	\begin{multicols}{3}
		\begin{enumerate}[label=\alph*)]
			\item $\log_2(50)$
			\item $4 \cdot \ln(4)$
			\item $a \log_b(c)$
		\end{enumerate}
	\end{multicols}
	
	\checkered[4cm]
	
	\begin{hint}{Exponentialfunktionen zur Basis e umschreiben}
		Jede Exponentialfunktion mit $f(x) = q^x$ lässt sich auch zur Basis $e$ schreiben:
		$$f(x) = q^x = \left(e^{\ln(q)}\right)^x = e^{\ln(q) \cdot x}= e^{k \cdot x} \text{ mit } k = \ln(q).$$
	\end{hint}
	
	\partnertask{Basiswechsel}
	
	Schreiben Sie in der jeweils anderen Form $q^x$ bzw. $e^{k \cdot x}$.
	
	\begin{multicols}{3}
		\begin{enumerate}[label=\alph*)]
			\item $1.5^x$
			\item $e^{0.5x}$
			\item $\frac{1}{4}^x$
			\item $e^{-2x}$
		\end{enumerate}
	\end{multicols}
	
	\checkered[4cm]
	
\end{document}